% --------------------------------------------------------------
% This is all preamble stuff that you don't have to worry about.
% Head down to where it says "Start here"
% --------------------------------------------------------------
 
\documentclass[12pt]{article}
 
\usepackage[margin=1in]{geometry} 
\usepackage{amsmath,amsthm,amssymb}
 
\newcommand{\N}{\mathbb{N}}
\newcommand{\Z}{\mathbb{Z}}
 
\newenvironment{theorem}[2][Theorem]{\begin{trivlist}
\item[\hskip \labelsep {\bfseries #1}\hskip \labelsep {\bfseries #2.}]}{\end{trivlist}}
\newenvironment{lemma}[2][Lemma]{\begin{trivlist}
\item[\hskip \labelsep {\bfseries #1}\hskip \labelsep {\bfseries #2.}]}{\end{trivlist}}
\newenvironment{exercise}[2][Exercise]{\begin{trivlist}
\item[\hskip \labelsep {\bfseries #1}\hskip \labelsep {\bfseries #2.}]}{\end{trivlist}}
\newenvironment{reflection}[2][Reflection]{\begin{trivlist}
\item[\hskip \labelsep {\bfseries #1}\hskip \labelsep {\bfseries #2.}]}{\end{trivlist}}
\newenvironment{proposition}[2][Proposition]{\begin{trivlist}
\item[\hskip \labelsep {\bfseries #1}\hskip \labelsep {\bfseries #2.}]}{\end{trivlist}}
\newenvironment{corollary}[2][Corollary]{\begin{trivlist}
\item[\hskip \labelsep {\bfseries #1}\hskip \labelsep {\bfseries #2.}]}{\end{trivlist}}
 
\begin{document}
 
% --------------------------------------------------------------
%                         Start here
% --------------------------------------------------------------
 
%\renewcommand{\qedsymbol}{\filledbox}
 
\title{Homework}%replace X with the appropriate number
\author{Amanda King\\ %replace with your name
Calculus 1 - Finding Limits Numerically} %if necessary, replace with your course title
 
\maketitle
 
\begin{exercise}{1} %You can use theorem, proposition, exercise, or reflection here.  Modify x.yz to be whatever number you are proving
$f(x)$ is given by formula: 
\end{exercise}
 
\begin{exercise}
1We are finding limits by using the "plug-in method" for numbers approaching 3 from the right :
%Note 1: The * tells LaTeX not to number the lines.  If you remove the *, be sure to remove it below, too.
%Note 2: Inside the align environment, you do not want to use $-signs.  The reason for this is that this is already a math environment. This is why we have to include \text{} around any text inside the align environment.
\begin{align*}
\lim_{x \to 3^+} (x^2+2) &= ?  \\
&  \text{Consider for }"x":  3.1,
               3.01,
               3.001,
               3.0001\text{} \\
& \lim(3.1^2+2)=11.61\\
& \lim(3.01^2+2)=11.0601\\
& \lim(3.001^2+2)=11.006001 \\
& \lim(3.0001^2+2)=11.0006
\end{align*}
\end{exercise}
 
\begin{proposition}{}
$\lim_{x \to 3^+} (x^2+2)$
Let $x=11$. 
\end{proposition}
 
\begin{proof}%Whatever you put in the square brackets will be the label for the block of text to follow in the proof environment.
Since all of the outcomes just keep getting closer to 11 one can assume that the limit as "x" approaches 3 from the right the limit will be equal to 11.
% * <arking4@g.coastal.edu> 2015-08-24T22:28:35.113Z:
%
% 
%
\end{proof}
 
% --------------------------------------------------------------
%     You don't have to mess with anything below this line.
% --------------------------------------------------------------
 
\end{document}